
\documentclass[a4paper,11pt]{article}
\usepackage{fontspec}
\setmainfont{Noto Sans}
\usepackage[french]{babel}
\usepackage{microtype}
\usepackage[top=9mm,bottom=8mm,left=11mm,right=11mm]{geometry}
\usepackage{xcolor}
\usepackage{tikz}
\usetikzlibrary{arrows.meta,calc,positioning}
\usepackage[most]{tcolorbox}
\usepackage{tabularx}
\usepackage{array}
\usepackage{ulem}
\pagestyle{empty}

\definecolor{Violet}{HTML}{7A5AA6}
\definecolor{VioletDark}{HTML}{563A7A}
\definecolor{VioletLight}{HTML}{F2ECF8}
\definecolor{VioletLine}{HTML}{B9A5D0}
\definecolor{Border}{HTML}{D3C4E2}
\definecolor{Ink}{HTML}{172A32}
\definecolor{Grey}{HTML}{64757D}
\definecolor{Red}{HTML}{D52F35}
\definecolor{RedLight}{HTML}{FCEBED}

\tcbset{
  enhanced,
  boxrule=.6pt,
  arc=2.8mm,
  left=3.7mm,right=3.7mm,top=1.75mm,bottom=1.75mm,
  boxsep=1mm,
  coltext=Ink
}

\newtcolorbox{violetbox}[1][]{
  colback=VioletLight,colframe=Border,#1
}
\newtcolorbox{warnbox}[1][]{
  colback=RedLight,colframe=Red,
  left=2.7mm,right=2.7mm,top=1.25mm,bottom=1.25mm,
  #1
}

\begin{document}

\begin{tikzpicture}[remember picture,overlay]
  % cadre extérieur : angles gauches carrés, droits légèrement arrondis
  \draw[draw=Border,line width=.9pt,rounded corners=1.7mm]
    ([xshift=1mm,yshift=-1mm]current page.north west)
    -- ([xshift=-1mm,yshift=-1mm]current page.north east)
    -- ([xshift=-1mm,yshift=1mm]current page.south east)
    -- ([xshift=1mm,yshift=1mm]current page.south west)
    -- cycle;
  % barre verticale
  \fill[Violet] ([xshift=1mm,yshift=-1mm]current page.north west)
    rectangle ([xshift=5.2mm,yshift=1mm]current page.south west);
\end{tikzpicture}

\vspace*{1mm}
\hspace*{6mm}
{\color{VioletDark}\fontsize{25}{27}\selectfont\bfseries CONDITIONNEL PRÉSENT}

\vspace{0.5mm}
\hspace*{6mm}
{\color{Grey}\large Le conditionnel présent}

\vspace{2.5mm}

\hspace*{6mm}
\begin{minipage}{0.91\textwidth}
\begin{violetbox}
\textbf{\color{VioletDark}CONSTRUCTION}\\[1mm]
\centering
{\Large \textbf{RADICAL DU FUTUR + TERMINAISONS DE L'IMPARFAIT}}\\[1mm]
\small
je \textbf{-ais}\quad tu \textbf{-ais}\quad il/elle \textbf{-ait}\quad
nous \textbf{-ions}\quad vous \textbf{-iez}\quad ils/elles \textbf{-aient}
\end{violetbox}

\vspace{2.5mm}

\begin{center}
\begin{tikzpicture}[>=Latex]
  \draw[Ink!65,line width=.8pt] (0,0) -- (15.1,0);
  \node[below=2mm] at (2.0,0) {\small PASSÉ};
  \node[below=2mm] at (7.5,0) {\small PRÉSENT};
  \node[below=2mm] at (13.5,0) {\small FUTUR};
  \draw[Violet,line width=2.5pt,->] (7.9,0) -- (12.9,0);
  \node[above=3mm,align=center,text width=6.2cm] at (10.4,0)
    {\small\bfseries possibilité, souhait, hypothèse, conseil};
  \draw[VioletDark,line width=1.1pt,->] (7.5,1.0) -- (7.5,0.18);
  \node[above=3.8mm] at (7.5,1.0) {\bfseries\color{VioletDark}MAINTENANT};
\end{tikzpicture}
\end{center}

\vspace{1.5mm}

\begin{violetbox}
\textbf{\color{VioletDark}À QUOI SERT-IL ?}\\[1mm]
\begin{itemize}
\setlength{\itemsep}{0.4mm}
\item exprimer une \textbf{hypothèse} ou une situation soumise à une condition ;
\item exprimer un \textbf{souhait}, un désir ou un rêve ;
\item donner un \textbf{conseil} ou formuler une demande avec atténuation ;
\item exprimer une information présentée avec \textbf{prudence} ou incertitude.
\end{itemize}
\end{violetbox}

\vspace{2mm}

\begin{violetbox}
\textbf{\color{VioletDark}EXEMPLES}\\[1mm]
\textbf{Si j'avais le temps, je voyagerais davantage.}\\
\textbf{J'aimerais visiter cette région.}\\
\textbf{Tu pourrais fermer la fenêtre ?}\\
\textbf{Il serait déjà parti.}
\end{violetbox}

\vspace{2mm}

\begin{minipage}{0.485\textwidth}
\begin{violetbox}[height=3.15cm]
\textbf{\color{VioletDark}AVEC « SI »}\\[1mm]
Une hypothèse au présent ou dans le futur se construit souvent avec :\\[1mm]
\centering
\textbf{SI + imparfait}\\
$\downarrow$\\
\textbf{conditionnel présent}\\[1mm]
\raggedright
\textit{Si j'avais plus de temps, je lirais davantage.}
\end{violetbox}
\end{minipage}
\hfill
\begin{minipage}{0.485\textwidth}
\begin{warnbox}[height=3.15cm]
\textbf{\color{Red}ATTENTION}\\[1mm]
Après \textbf{si}, on ne met normalement pas le conditionnel dans ce type de phrase.\\[1mm]
\textbf{Si j'avais...}\\
\textcolor{Grey}{et non : \sout{Si j'aurais...}}
\end{warnbox}
\end{minipage}

\vspace{2mm}

\begin{violetbox}
\textbf{\color{VioletDark}À RETENIR}\\[1mm]
Le conditionnel présent sert à présenter une action comme \textbf{possible, souhaitée,
hypothétique ou atténuée}. Il ne désigne pas simplement un « futur dans le passé » :
sa valeur dépend du contexte.
\end{violetbox}

\vspace{1.5mm}

\centering
{\small\color{Grey}\textbf{Repères fréquents :} si, dans ce cas, à ta place, volontiers, peut-être, probablement, j'aimerais, je voudrais, je pourrais...}

\end{minipage}

\end{document}
